\chapter{Basis-set / ECP recommendation (menu 4)}
\label{ch:menu4}
\menulabel{4}

\section{Purpose}

Look up the basis-set and ECP recommendation for an element set, valence and
target without generating a file -- the same knowledge that \menu{3} uses.

\section{What you need}

Element list, charge, multiplicity, f-block metal valence, targets (the same
profile as \menu{3}).

\section{Walkthrough}

\lstinputlisting[style=fbkcode, firstline=16, lastline=36,
  title={The menu-4 example session (Pu, Cl)}]
  {examples/expected/m4_basis.txt}

\section{What you get}

Recommendation tiers (each with its matching auxiliary basis and rationale),
boundary cautions, and not-applicable tiers with the hard refusal conditions --
everything with provenance. The example shows the three shapes of an answer:

\begin{itemize}
  \item \textbf{Tiers}: SARC-DKH-TZVPP first for an actinide (SARC covers
    Ac--Lr with its own SARC/J fitting basis), ANO-RCC as the
    magnetic-property-chain generalist, Stuttgart pseudopotentials through
    ORCA's built-in keywords.
  \item \textbf{Boundary cautions}: the standard def2 family stops at Rn and
    does not cover the actinides -- do not pick it for an actinide.
  \item \textbf{Not applicable, with reasons}: the 5f-in-core pseudopotential
    family applies to the trivalent state only (its own calibration: AnCl$_3$
    MAD 2.9 kcal/mol against small-core pseudopotentials; the tetravalent case
    does not meet the bar), and it is structurally blind to everything that
    needs an f channel with $l_{\max} = 2$: f--f transitions, crystal-field
    splittings, 5f multiplets, every SOC-dependent property. Its refusal list
    also states the source paper's own prohibition (ground state
    $5f^{n-1}6d^1$).
\end{itemize}

\section{How to read the numbers}

\begin{readbox}
The not-applicable block is the part to read carefully: it is the program
declining to recommend, with the quantified reason attached. A tier that is
``not applicable'' for one target may be exactly right for another (the
5f-in-core family is legitimate for closed-shell chemistry with no f physics);
read the condition, not the verdict alone.
\end{readbox}

\section{Worked example}

Run the session above; then change the valence from 3 to 4 and the targets to
\code{excited} and compare the blocks. The recommendation changes because the
rules match on (elements, valence, targets) -- there is no hidden state.

\section{Caveats, refusals and related sections}

\begin{refusalbox}
\begin{itemize}
  \item This is a lookup, not a calculation: no file is written.
  \item ECP parameters from an external source (a Basis Set Exchange download,
    say) must be inlined into the ORCA input and checked by hand -- the
    program says so where that route is suggested.
  \item Where the licence or availability of an external basis was not
    verified, the entry says ``to verify'' instead of asserting it.
\end{itemize}
\end{refusalbox}

Related: \menu{3} (generation), Chapter~\ref{ch:criteria} (the underlying
rules and their sources).
